\documentclass{article}

\usepackage{algorithm}
\usepackage{algpseudocode}
\usepackage{appendix}
\usepackage{amsmath}
\usepackage{amssymb}
\usepackage[UTF8]{ctex}
\usepackage{fancyhdr}
\usepackage{float}
\usepackage{graphicx}
\usepackage{listings}
\usepackage{multirow}
\usepackage{tabularx}
\usepackage{newclude}
\usepackage{subfig}
\usepackage{xcolor}

\lstset{
    basicstyle          =   \sffamily,          % 基本代码风格
    keywordstyle        =   \bfseries,          % 关键字风格
    commentstyle        =   \rmfamily\itshape,  % 注释的风格，斜体
    stringstyle         =   \ttfamily,  % 字符串风格
    flexiblecolumns,                % 别问为什么，加上这个
    numbers             =   left,   % 行号的位置在左边
    showspaces          =   false,  % 是否显示空格，显示了有点乱，所以不现实了
    numberstyle         =   \zihao{-5}\ttfamily,    % 行号的样式，小五号，tt等宽字体
    showstringspaces    =   false,
    captionpos          =   t,      % 这段代码的名字所呈现的位置，t指的是top上面
    frame               =   lrtb,   % 显示边框
}

\lstdefinestyle{Python}{
    language        =   Python, % 语言选Python
    basicstyle      =   \zihao{-5}\ttfamily,
    numberstyle     =   \zihao{-5}\ttfamily,
    keywordstyle    =   \color{blue},
    keywordstyle    =   [2] \color{teal},
    stringstyle     =   \color{magenta},
    commentstyle    =   \color{red}\ttfamily,
    breaklines      =   true,   % 自动换行，建议不要写太长的行
    columns         =   fixed,  % 如果不加这一句，字间距就不固定，很丑，必须加
    basewidth       =   0.5em,
}

\lstdefinestyle{Bash}{
    language        =   Bash, % 语言选Python
    basicstyle      =   \zihao{-5}\ttfamily,
    numberstyle     =   \zihao{-5}\ttfamily,
    keywordstyle    =   \color{blue},
    keywordstyle    =   [2] \color{teal},
    stringstyle     =   \color{magenta},
    commentstyle    =   \color{red}\ttfamily,
    breaklines      =   true,   % 自动换行，建议不要写太长的行
    columns         =   fixed,  % 如果不加这一句，字间距就不固定，很丑，必须加
    basewidth       =   0.5em,
}

\title{计算机视觉——第二次作业}
\author{Koorye}
\date{2024年4月8日}

\pagestyle{fancy}
\fancyhead[L]{Local Feature Matching}
\fancyhead[R]{Koorye}

\begin{document}
% \maketitle

\thispagestyle{empty}

\begin{figure}[t]
    \centering
    \includegraphics[width=0.8\textwidth]{images/logo.png}
\end{figure}

\vspace*{\fill}
    \begin{center}
        \huge\textbf{高级计算机视觉——第二次作业}
    \end{center}
\vspace*{\fill}

\begin{table}[b]
    \centering
    \large
    \begin{tabular}{ll}
    \textbf{课程:} & 高级计算机视觉 \\
    \textbf{姓名:} & Koorye \\
    \textbf{学号:} & xxxxxx \\
    \textbf{时间:} & 2024年4月9日 \\
    \end{tabular}
\end{table}

\newpage
\tableofcontents

\newpage
\include*{chapters/introduction}

\include*{chapters/principle}

\include*{chapters/steps}

\include*{chapters/results}

\include*{chapters/conclusion}

\newpage
\include*{chapters/appendix}

\end{document}